% from ega1-3-fr.tex
\oldpage[I]{108}
précède~; mais alors les produits $X_i\times_S Y_i$ existent aussi
(\hyperref[I.3.2.5-fr]{3.2.5}), donc il en est de même de $X\times_S Y$ par
(\hyperref[I.3.2.6.4-fr]{3.2.6.4}).
\end{env}

\begin{corollary}[3.2.7]
\label{I.3.2.7-fr}
Soient $Z=X\times_S Y$ le produit de deux $S$-préschémas, $p$, $q$ les
projections de $Z$ dans $X$ et $Y$, $\varphi$ (resp. $\psi$) le morphisme
structural de $X$ (resp. $Y$). Soient $S'$ une partie ouverte de $S$, $U$
(resp. $V$) une partie ouverte de $X$ (resp. $Y$) contenue dans
$\varphi^{-1}(S')$ (resp. $\psi^{-1}(S')$). Alors le produit
$U\times_{S'}V$ s'identifie canoniquement au préschéma induit par $Z$ sur
$p^{-1}(U)\cap q^{-1}(V)$ (considéré comme $S'$-préschéma). En outre, si
$f:T\to X$, $g:T\to Y$ sont des $S$-morphismes tels que $f(T)\subset U$,
$g(T)\subset V$, le $S'$-morphisme $(f,g)_{S'}$ s'identifie à la restriction
de $(f,g)_S$ à $p^{-1}(U)\cap q^{-1}(V)$.
\end{corollary}

Cela résulte de (\hyperref[I.3.2.5-fr]{3.2.5}) et
(\hyperref[I.3.2.6.1-fr]{3.2.6.1}).

\begin{env}[3.2.8]
\label{I.3.2.8-fr}
Soient $(X_\alpha)$, $(Y_\lambda)$ deux familles de $S$-préschémas, $X$
(resp. $Y$) la somme de la famille $(X_\alpha)$ (resp. $(Y_\lambda)$)
(\hyperref[subsection:I.3.1-fr]{3.1}). Alors $X\times_S Y$ s'identifie à la
\emph{somme} de la famille $(X_\alpha\times_S Y_\lambda)$~; cela résulte
aussitôt de (\hyperref[I.3.2.6.3-fr]{3.2.6.3}).
\end{env}

\subsection{Propriétés formelles du produit~; changement de préschéma de base.}
\label{subsection:I.3.3-fr}

\begin{env}[3.3.1]
\label{I.3.3.1-fr}
Le lecteur remarquera que toutes les propriétés énoncées dans cette section,
sauf (\hyperref[I.3.3.13-fr]{3.3.13}) et
(\hyperref[I.3.3.15-fr]{3.3.15}), sont valables sans modification dans
\emph{toute catégorie}, chaque fois que les produits qui interviennent dans
les énoncés \emph{existent} (car il est clair que les notions de $S$-objet et
de $S$-morphisme peuvent se définir exactement comme dans
(\hyperref[subsection:I.2.5-fr]{2.5}) pour tout objet $S$ de la catégorie).
\end{env}

\begin{env}[3.3.2]
\label{I.3.3.2-fr}
En premier lieu, $X\times_S Y$ est un \emph{bifoncteur covariant} en $X$ et
$Y$ dans la catégorie des $S$-préschémas~: il suffit en effet de remarquer
que le diagramme
\[
  \xymatrix{
    X\times Y\ar[r]^{f\times 1}\ar[d] &
    X'\times Y\ar[r]^{f'\times 1}\ar[d] &
    X''\times Y\ar[d]\\
    X\ar[r]_f & X'\ar[r]_{f'} & X''
  }
\]
est commutatif.
\end{env}

\begin{proposition}[3.3.3]
\label{I.3.3.3-fr}
Pour tout $S$-préschéma $X$, la première (resp. seconde) projection de
$X\times_S S$ (resp. $S\times_S X$) est un isomorphisme fonctoriel de
$X\times_S S$ (resp. $S\times_S X$) sur $X$, dont l'isomorphisme réciproque
est $(1_X,\varphi)_S$ (resp. $(\varphi,1_X)_S$), en désignant par $\varphi$
le morphisme structural $X\to S$~; on peut donc écrire, à un isomorphisme
canonique près
\[
  X\times_S S=S\times_S X=X.
\]
\end{proposition}

Il suffit de prouver que le triplet $(X,1_X,\varphi)$ est un produit de $X$
et $S$. Or, si $T$ est un $S$-préschéma, le seul $S$-morphisme de $T$ dans
$S$ est nécessairement le morphisme structural $\psi:T\to S$. Si $f$ est
un $S$-morphisme de $T$ dans $X$, on a nécessairement
$\psi=\varphi\circ f$, d'où notre assertion.

\begin{corollary}[3.3.4]
\label{I.3.3.4-fr}
Soient $X$, $Y$ deux $S$-préschémas, $\varphi:X\to S$, $\psi:Y\to S$ leurs
morphismes structuraux. Si on identifie canoniquement $X$ à $X\times_S S$
et $Y$ à $S\times_S Y$, les projections $X\times_S Y\to X$ et
$X\times_S Y\to Y$ s'identifient respectivement à $1_X\times\psi$ et
$\varphi\times 1_Y$.
\end{corollary}

La vérification est immédiate et laissée au lecteur.

\begin{env}[3.3.5]
\label{I.3.3.5-fr}
On peut définir de la même manière que dans
(\hyperref[subsection:I.3.2-fr]{3.2}) le produit d'un
